% !PRINTED 689
\weberpage{689}
%<<CONT
und wenn $y_k, y_k'$ einem imaginären Paare angehören,
\weq{9}{z_k = \log y_k y_k',}
wodurch die Definition völlig eindeutig ist.

Nun bestimmen wir die Variablen $\xi_1, \xi_2, \ldots, \xi_{\nu-1}, \xi_\nu$ aus den linearen Gleichungen
\weq{10}{\begin{array}{l}
\xi_1 l_{1,1} + \cdots + \xi_{\nu-1} l_{\nu-1,1} + \delta_1 \xi_\nu = z_1\\
\xi_1 l_{1,2} + \cdots + \xi_{\nu-1} l_{\nu-1,2} + \delta_2 \xi_\nu = z_2\\
. \; . \; . \; . \; . \; . \; . \; . \; . \; . \; . \; . \; . \; .\\
\xi_1 l_{1,\nu} + \cdots + \xi_{\nu-1} l_{\nu-1,\nu} + \delta_\nu \xi_\nu = z_\nu,
\end{array}}
deren Determinante nicht verschwindet.

Giebt man den Variablen $x_i$ ganzzahlige Werthe, so gehen nach \S.\,188, (2), (3) die $\xi_1, \xi_2, \ldots, \xi_{\nu-1}$ in das Exponentensystem der Zahl $\alpha$ über, und es ergiebt sich durch Addition der Gleichungen (10) [\S.\,185, (3)]:
\weq{11}{\xi_\nu = \frac{1}{n} \log N_a(\alpha).}

Wir gehen jetzt von dem Punkte $(x)$ zu dem Punkte $(x')$ über, indem wir $x_i' = t^{\frac{1}{n}} x_i$ setzen. Dadurch geht $y_i$ in
\[ y_i' = t^{\frac{1}{n}} y_i \]
über, und $z_i$ in
\[ z_i' = z_i + \frac{1}{n} \delta_i \log t. \]

Wenn wir also $z_i'$ an Stelle von $z_i$ setzen, und die dann aus (10) sich ergebenden Werthe der $\xi_1, \xi_2, \ldots, \xi_\nu$ mit $\xi_1', \xi_2', \ldots, \xi_\nu'$ bezeichnen, so ist
\weq{12}{\xi_1' = \xi_1, \xi_2' = \xi_2, \ldots, \xi_{\nu-1}' = \xi_{\nu-1}, \xi_\nu' = \xi_\nu + \frac{1}{n} \log t.}

Nunmehr wollen wir das Gebiet $S'$ dadurch abgrenzen, dass
\weq{13}{\begin{array}{c}
0 \leqq \xi_1' < 1, \ldots, 0 \leqq \xi_{\nu-1}' < 1, \\
\xi_\nu' < \frac{1}{n} \log t
\end{array}}
sein soll. Dadurch erreichen wir, dass die in dem Gebiete $S'$ liegenden Gitterpunkte $(x')$ nur reducirte ganze Zahlen $\alpha$ darstellen, und alle und nur solche, deren absolute Norm kleiner als $t$ ist.

Unter den reducirten Zahlen sind aber immer je $w$ mit einander associirt, wenn $w$ die Anzahl der in $\Omega$ enthaltenen
%CONT>>
